% 原创教学示意；同一发送状态被两条消息路径读取，反向贡献相加。
\begin{tikzpicture}[x={\dimexpr\linewidth/169\relax},y=1mm,
  font=\figurefont\footnotesize,text=NNInk]
  \node[nn operator,minimum width=19mm,minimum height=14mm]
    (sender) at (12,0) {$\bm h_j^{(l)}$};
  \draw[nn flow,draw=NNHidden] (sender.east)--(30,0);
  \foreach \idx/\yy in {i/18,k/-18}{
    \node[nn operator,minimum width=24mm,minimum height=14mm]
      (message-\idx) at (52,\yy) {$M^{(l)}$};
    \node[nn operator,minimum width=30mm,minimum height=14mm]
      (update-\idx) at (101,\yy) {聚合与更新};
    \node[nn operator,minimum width=22mm,minimum height=14mm]
      (state-\idx) at (154,\yy) {$\bm h_{\idx}^{(l+1)}$};
    \draw[nn flow,draw=NNHidden] (30,0)--([yshift=3mm]message-\idx.west);
    \draw[nn flow,draw=NNHidden]
      ([yshift=3mm]message-\idx.east)--([yshift=3mm]update-\idx.west);
    \draw[nn flow,draw=NNHidden]
      ([yshift=3mm]update-\idx.east)--([yshift=3mm]state-\idx.west);
    \draw[nn gradient]
      ([yshift=-3mm]state-\idx.west)--([yshift=-3mm]update-\idx.east);
    \draw[nn gradient]
      ([yshift=-3mm]update-\idx.west)--([yshift=-3mm]message-\idx.east);
    \node[text=NNGradient,anchor=north] at (78,\yy-7)
      {$\bm\delta_{\idx\leftarrow j}^{(l)}$};
    \node[text=NNGradient,anchor=north] at (130,\yy-7)
      {$\bm g_{\idx}^{(l+1)}$};
  }
  \draw[nn gradient] ([yshift=-3mm]message-i.west)
    to[out=180,in=0] ([yshift=4mm]sender.east);
  \draw[nn gradient] ([yshift=-3mm]message-k.west)
    to[out=180,in=0] ([yshift=-4mm]sender.east);
  \node[text=NNGradient,anchor=north] at (12,-11) {$\bm g_j^{(l)}$};
\end{tikzpicture}
